% Appendix B --- Verification: the inventory behind section 1.1 (Q-0414, round 1; the author's
% decision (c) of 2026-10-11; article-kit WRITING.md section 4 and PUBLICATION-TEMPLATE.md B.4).
% Section 1.1 states the trust base: what is machine-checked, what is assumed, what is neither.
% This appendix adds detail to that statement and holds no caveat, gap or assumption that 1.1
% does not state: where the formal development is, the versions, the two commands, the Lean names
% of the seven cited facts, and the table of the cited facts per statement. The statements
% covered and Lean's own `#print axioms` output are section S4 of the supplement
% (supp-verification.tex) since the split of Q-0448.
%
% Moved here from 01-introduction.tex by records/affine/edits_q0414_verification_appendix.py,
% which copies the table of identifiers and the stamped block byte for byte and checks that they
% are identical to the blocks removed from section 1.1. The stamped block's content is untouched:
% no build is available in the worktree of this item, so `linkage check --fresh-axioms` was not
% run, and no line was added (the block has no line for the rate's declarations; adding one needs
% a build and a new stamp).
%
% The region between the two "trust base" comment lines is read by `linkage check` (LINKAGE.md
% rule 13), as the region of section 1.1 is: a paper may have several regions, in one file or
% several. The sentence that lists the statements by \ref is what the status comparison reads;
% the names of the cited facts are what the ledger comparison reads, against the axioms that the
% listed statements' declarations print; the block is compared with the #guard_msgs pins of
% Formalization/CIAxiomGuard.lean.

\section{Verification}\label{app:verification}
% trust base: begin

% Q-0428 (the presentation review of round 1, item 19): the entry in plain words, with the three
% counts together. That no declaration of the list depends on the three other axioms is said
% under "The report for every declaration", from the report of the 141 declarations.
% Q-0431 (the author's section review of 2026-10-12, records/affine/REVIEW-author-paper-iso.md:
% "we focus on what is not what isn't"): the appendix read once against that principle. Four
% sentences whose content was what is not the case are turned or removed: "It adds no
% assumption to what that subsection states" (every assumption is stated there); "The four
% statements of its second list have no item" (the paragraph already says which statements the
% list covers); "it has no request for a definition. So it does not print the report of the
% last paragraph" (of the guard: the file of the report is named in the paragraph before); "It
% has no declaration for 4 <= w_2" (that inequality is the computed one). Kept, because a
% reader needs them to know what is checked: "no file of this development uses it" (the
% earlier papers' development in the build configuration), "none depends on any other axiom",
% "on which no declaration of this paper depends", and that the repository is private.
This appendix says where the formal development behind \S\ref{sec:trust-base} is, in which
versions, under which names it holds the cited facts, and which cited facts each statement
rests on. \suppl{} has the declarations statement by statement, with Lean's report
(\S\ref{supp:verification}).
A \emph{declaration} is a named definition or proved statement of the formal
development. An \emph{axiom}, in Lean's report, is a declaration taken without proof, and the
report on a declaration lists every axiom that its proof uses. The seven cited facts of this
paper are nine such axioms, C3 being three. The guard of the repository admits twelve, because
the formal development also states the three cited facts of the companion paper.
% Q-0452: "because it also serves the companion paper, which has three cited facts of its own"
% until the isotropic paper's Lean was separated from the affine paper's: the guard of the public
% repository lists no declaration of the companion paper, and admits twelve axioms because
% Formalization/AffineHemigroup/Interfaces.lean states all the cited facts of chapter 2 and
% blueprint/trust-boundary.txt lists them. True of the development repository as well.

% Q-0414, round 3 (the register review of 2026-10-11b, flags 13-15, 25, 26, 36, 37; the scope
% audit's findings 11 and 12): the list of statements is a table with the Lean declarations,
% generated from the blueprint's \lean tags by records/affine/gen-verification-table.py (`check`
% compares); the pass criterion is one sentence; the paths are from the repository root; nine
% axioms for seven cited facts; the formal examples are named; the build configuration is said as
% Formalization/lakefile.toml has it (the public export of the cone paper, tag v0.1, which
% contains the line paper's development).
\paragraph{Where the formal development is.} The formal development, its ledger of the cited
facts with the page-verified citation of each, and the blueprint, the sources in which the
statements and their proofs were developed, are in the repository accompanying this paper (see
the code availability statement at the end). They are the directories \texttt{Formalization/}
and \texttt{blueprint/} of that repository. The repository is
\url{https://github.com/danielfagerstrom/spatial-hemigroup-scale-space-isotropic},
and this version of the paper is its tag \texttt{v0.1}. It is a public verification export: it is
taken from the development repository, which is private, and it is
deposited with this version of the paper under the DOI printed on its first page. It holds
the formal development, the ledger, the blueprint sources of the two chapters behind this
paper, and the scripts of the figures and of Table~\ref{tab:matern}, each at the path that this
paper or the supplement prints.
% The release iso-v0.1 (records/affine/freeze_iso_release.py); the draft form is in git, and its
% sentences are quoted in the comment of the code availability statement (declarations.tex).
% Q-0417 (the freeze for external review, article-kit PROCESS.md section 9): the repository and
% the revision, named by a tag because a commit cannot print its own hash; the dispatching
% session puts the tag iso-review-1 on the merge commit of this item. The repository is private,
% so its name is printed as a name, not as a URL (the scope audit's row D3 and the register
% review's flag 15 asked that a reviewer of the frozen text can tell what this appendix points
% to). The same three facts are in the code availability statement.
% TODO(release): the URL and tag of the public verification export, as Paper V names its export
% (the repository itself is private), and the review-build sentences above replaced by them.


\paragraph{Versions and commands.} The versions are fixed in the repository: Lean~4
\texttt{v4.31.0}, Mathlib \texttt{v4.31.0}, and the shared library \texttt{ScaleSpaceCore} at its
tag \texttt{v0.4.0}. The build configuration also names the formal development of the two earlier
papers, from the public export of the cone paper at its tag \texttt{v0.1}, and no file of this
development uses it. \texttt{lake build} in \texttt{Formalization/} reproduces the verification:
it checks every proof. On request Lean reports, for a declaration, the assumptions it depends on,
which it calls its axioms. The criterion for this paper is that every axiom reported for a
declaration of the list of \S\ref{supp:verification} is one of Lean's three logical axioms,
\texttt{propext},
\texttt{Classical.choice} and \texttt{Quot.sound}, or one of the nine names below.
% Q-0424 (the external review, round 1; records/affine/reviews/
% VERIFY-round1-claims-and-trust-base.md, item 6(b)). Until now the criterion printed here was
% the repository guard's, which admits twelve axioms, three of them the companion paper's, and
% the text said that `lake env lean CIAxiomGuard.lean` prints the report for every named
% declaration. It does not: of the 141 declarations the guard file has a line for 99, pins 31 of
% those under #guard_msgs, which prints nothing when the pin holds, and has no line for the 42
% definitions; the eight of the printed block are among the pinned. The paragraphs below state
% the criterion of this paper (nine names), name the file that holds Lean's report for all 141
% declarations and the script that makes and checks it, and say what the guard is. A guard of
% this paper alone, in CI, is another item of the response plan (R2).

\paragraph{The names of the cited facts.} The nine names are the seven cited facts of this
paper, C3 being three. Table~\ref{tab:trust-base} says which of them each statement rests on.
They are in the namespace \texttt{AffineHemigroup}:
\begin{center}
\small
\begin{tabular}{@{}ll@{\qquad}ll@{}}
C1 & \texttt{density\_of\_integrable\_charFun} & C3(iii) & \texttt{levy\_khintchine\_converse} \\
C2 & \texttt{sato\_polar\_criterion} & C4 & \texttt{selfDecomposable\_infinitelyDivisible} \\
C3(i) & \texttt{levy\_khintchine\_exists} & C5 & \texttt{wolfe\_symmetric\_unimodal} \\
C3(ii) & \texttt{levy\_khintchine\_unique} & C6 & \texttt{schoenberg\_radial\_bernstein} \\
 & & C7 & \texttt{karamata\_tauberian\_zero}
\end{tabular}
\end{center}

% Q-0428: Table tab:trust-base, moved here from section 1.1 with the statements of Q-0425 and
% Q-0429 added. The status of each statement is in section 1.1 and is not repeated in the rows.
% At the merge of the Lean of Q-0429 the rows of the six statements formalized since the first
% review build were compared with Lean's report
% (records/affine/axioms-iso-all-declarations.py, `TABLE1`, which the generated guard
% Formalization/AffineHemigroup/IsoAxiomGuard.lean enforces at every build): C2, C3(iii), C5 and
% C7 for the rate through the profile; C2, C3(i) to (iii) and C4 for
% prop:profile-one-dimension-down and for thm:separating-family-critical-values; none for
% prop:isotropic-exponent-profile, lem:threshold-asymptotic and lem:shadow-polar-profile, which
% are under "every other statement". No row changed. The row of the dictionary holds one
% statement with a formal counterpart, prop:profile-one-dimension-down, beside two without; the
% cited facts of those two are the ones their printed proofs use, as the caption says.
% Q-0448, the joining step (the Lean of Q-0447 merged): the row of the stages compared with
% Lean's report, declaration by declaration. The twelve declarations of clause (1) and the six
% of clause (2) name no cited fact; of clause (3), isotropic_stage_triplet,
% isotropic_stage_absolutelyContinuous and isotropic_stage_no_point_mass name C3(iii) (items
% (a) and (b)), isotropic_stage_smooth_density names C1 (item (c)), and
% isotropic_stage_dichotomy, the closing sentence, names both. The row stands as it was written
% from the printed proof; its entry in `TABLE1` is C1 and C3(iii).
\begin{table}[!htb]
\centering
\small
\begin{tabular}{@{}p{0.46\textwidth}p{0.5\textwidth}@{}}
\toprule
statement & cited facts its proof rests on \\
\midrule
Thm.~\ref{thm:isotropic-classification} and Cor.~\ref{cor:isotropic-classification-operators},
  the classification & C3(ii) and (iii), on the line \\
Cor.~\ref{cor:isotropic-scale-zero-kernels}, the main result as one statement & C2, C3, C4 \\
Prop.~\ref{prop:isotropic-polar-reading}, the polar form of $\mathcal{A}_d$ & C2, C3, C4 \\
Thm.~\ref{thm:origin-boundedness-d}, the threshold & C1, in clauses (2) and (5) \\
Prop.~\ref{prop:isotropic-radial-monotone}, the radially nonincreasing density & C5; C3(ii) and
  (iii) in clause (1); C2 and C3(iii) in clause (2) \\
Prop.~\ref{prop:origin-lower-bound-d}, the lower bound & C3(iii) \\
Prop.~\ref{prop:origin-upper-bound-d}, the upper bound & none for the bounds of clauses (1)
  and (2); C3(iii) for clause (3); C2, C3(iii) and C5 for the existence of the density \\
Lemma~\ref{prop:origin-rate-d} of \suppl, the rate through the transform & C7 for the limits; C2,
  C3(iii) and C5 for the existence of the density \\
Prop.~\ref{prop:origin-rate-profile-d}, the rate through the profile & C7 for the limits; C2,
  C3(iii) and C5 for the existence of the density \\
Prop.~\ref{prop:isotropic-stages}, the stages & C3(iii) in clause (3), items (a) and (b); C1
  in its item (c); none in clauses (1) and (2) \\
Prop.~\ref{prop:radial-dictionary}, the dictionary, Prop.~\ref{prop:profile-one-dimension-down},
  Cor.~\ref{cor:dictionary-filtration} and Prop.~\ref{prop:step-exponent-dimensions} & C2, C3,
  C4 \\
Thm.~\ref{thm:separating-family-critical-values}, the separating family & C2, C3, C4, in
  clauses (1) and (4) \\
Prop.~\ref{prop:separating-family-mixtures}, the separating family as laws & C2, C3, C4, in
  clause (3) \\
Lemma~\ref{lem:filtration-top} of \suppl, the top of the filtration,
  $\mathcal{A}_\infty = \mathcal{S}$ & C2, C3, C4, C6 \\
every other statement & none \\
\bottomrule
\end{tabular}
\caption{The statements of the paper and of the supplement and the cited facts that their
proofs rest on. For a statement with a formal counterpart these are
the assumptions that Lean reports for its declarations, and for the statements of the second
list of \S\ref{sec:trust-base} the cited facts that the printed proof uses. ``None'' means none of C1
to C7: the statement follows from the definitions and the logical foundations of the proof
assistant alone, or, where there is no formal counterpart, its printed proof cites no fact.}
\label{tab:trust-base}
\end{table}


% Q-0428, the fixes round: the account statement by statement, moved here from section 1.1,
% which keeps the caveat in two sentences (the register review of the front, flag 28;
% article-kit WRITING.md section 4: the appendix adds detail to what 1.1 says).
\paragraph{The statements restated from the two earlier papers.} The printed proofs cite the
statements restated in \S\ref{sec:line}, and the formal proofs assume none of them. So in two
places the formal proof derives from the cited facts what the printed proof cites from an
earlier paper: the uniqueness of the \Levy{} measure in the classification on the line, from
C3, and the identification $\mathcal{A}_\infty = \mathcal{S}$, from C6 with C2, C3 and C4. In
their own papers the restated statements stand as follows. Of those of the line paper, the
passage from operators to kernels, the classification and the two facts used in proofs are
verified formally there, the classification on two cited facts, and that the admissible
exponents are those of the symmetric self-decomposable laws is cited from the literature
(\sscite{fagerstrom2026spatial}, \S1.1). Of those of the cone paper, the bridge on exponents,
the extreme rays and the threshold on the line are verified formally there, on cited facts that
it lists, and the dimension filtration is proved in its text (\sscite{fagerstrom2026cone},
\S1.1).

% trust base: end
